%% ============================================================
% f2012_sichuan_arts_q21_locus.tex
%% 2012 四川卷（文科）第 21 题 动点 M 与定点 A(-1,0)、B(1,0)
%% 来源 PDF：普通高考/2012/2012四川文.pdf 第 2 页（印面 774）第 2 栏
%% 坐标取自原页矢量层（bp），原点为图区左上角 (630.0,500.0)，y 向下。
%% ============================================================
\begin{tikzpicture}[x=1bp,y=-1bp,
  every node/.style={inner sep=0,font=\fontsize{10.5pt}{12.6pt}\selectfont},
  line cap=butt, line join=miter]
  \draw[line width=0.3985bp] (7.656,85.133) -- (135.832,85.133);
  \fill[black] (139.418,85.133) .. controls (138.355,84.934) and (136.629,84.336) .. (135.434,83.641) -- (135.434,86.629) .. controls (136.629,85.93) and (138.355,85.332) .. (139.418,85.133);
  \fill[black] (64.328,5.793) .. controls (64.129,6.855) and (63.531,8.582) .. (62.832,9.777) -- (65.82,9.777) .. controls (65.125,8.582) and (64.527,6.855) .. (64.328,5.793);
  \draw[line width=0.3985bp] (35.992,85.133) -- (106.832,11.461);
  \draw[line width=0.3985bp] (92.664,85.133) -- (106.832,11.461);
  \node[anchor=base west] at (55.37,9.16) {$y$};
  \node[anchor=base west] at (110.89,12.05) {$M$};
  \node[anchor=base west] at (136.71,95.51) {$x$};
  \node[anchor=base west] at (89.79,97.76) {$B$};
  \node[anchor=base west] at (29.3,98.05) {$A$};
  \node[anchor=base west] at (51.74,99.59) {$O$};
\end{tikzpicture}
